\chapter{Glosario y acrónimos}
\label{anexo:glosario}

Este glosario recoge las abreviaturas, siglas y términos técnicos que
aparecen en la memoria, con una breve definición contextualizada al
trabajo.

\section*{Acrónimos técnicos}
\addcontentsline{toc}{section}{Acrónimos técnicos}

\begin{description}
  \item[ADC] Adenocarcinoma. Subtipo histológico de NSCLC caracterizado
  por su origen glandular y por la expresión de marcadores del programa
  alveolar tipo II.
  \item[API] Application Programming Interface. En este trabajo se
  emplea la API pública de NCBI GEO para descargar cohortes y la API de
  Groq para consultar al modelo Llama.
  \item[AUC] Area Under the (ROC) Curve. Métrica que resume la
  capacidad de un clasificador binario para ordenar correctamente pares
  positivo/negativo. Rango $[0,1]$; 0,5 equivale al azar y 1,0 al
  clasificador perfecto.
  \item[Bal.\,acc.] Balanced accuracy. Media aritmética de sensibilidad
  y especificidad; robusta frente al desbalance de clases.
  \item[cDNA] DNA complementario. Producto de la retrotranscripción de
  ARN a ADN usado como sustrato para secuenciación en RNA-Seq.
  \item[ComBat] Herramienta estándar de corrección de efecto lote basada
  en modelos bayesianos empíricos. No se usa en este trabajo por
  incompatibilidad conceptual con LODO.
  \item[CV] Cross-validation. Validación cruzada. En este trabajo se
  reserva a estudios previos de calibración de hiperparámetros; la
  validación externa se realiza con LODO.
  \item[$d$ (de Cohen)] Tamaño de efecto estandarizado para diferencias
  de medias. Un $|d|>0{,}8$ se considera un efecto grande. Se usa como
  métrica principal del criterio de consenso multi-cohorte.
  \item[DEG] Differentially Expressed Gene. Gen cuya expresión difiere
  significativamente entre dos grupos experimentales.
  \item[FDR] False Discovery Rate. Corrección estadística para múltiples
  comparaciones (Benjamini-Hochberg en este trabajo).
  \item[GEO] Gene Expression Omnibus. Repositorio público del NCBI que
  aloja las cohortes transcriptómicas empleadas en el trabajo.
  \item[GEOparse] Librería Python que consulta la API de GEO y devuelve
  matrices de expresión y metadatos en formato \texttt{pandas}.
  \item[GPL570] Identificador de plataforma en GEO correspondiente al
  chip Affymetrix Human Genome U133 Plus 2.0, la plataforma dominante
  en este trabajo.
  \item[HUGO] Human Genome Organisation. Autoridad de nomenclatura
  génica cuyos símbolos oficiales se emplean como identificadores de
  gen en la firma.
  \item[IARC] International Agency for Research on Cancer, agencia
  especializada de la OMS que publica estadísticas globales de cáncer.
  \item[IHC] Inmunohistoquímica. Técnica clásica de diagnóstico
  histopatológico que emplea anticuerpos para detectar la expresión de
  proteínas en secciones de tejido.
  \item[LASSO] Least Absolute Shrinkage and Selection Operator.
  Regresión con regularización $\ell_1$ que fuerza coeficientes a cero,
  induciendo selección de variables. En este trabajo se usa como
  clasificador principal para la firma.
  \item[LCNE] Large Cell Neuroendocrine Carcinoma. Subtipo de tumor
  neuroendocrino de pulmón, fuera del alcance de entrenamiento del
  framework.
  \item[LLM] Large Language Model. Modelo de lenguaje de gran tamaño
  (Llama~3.3-70b en este trabajo) empleado para la curación clínica de
  metadatos y como asistente conversacional en la interfaz.
  \item[LODO] Leave-One-Dataset-Out. Esquema de validación externa en el
  que cada cohorte se retira una vez del entrenamiento y se usa como
  test; el proceso se itera sobre todas las cohortes evaluables.
  \item[NSCLC] Non-Small Cell Lung Cancer. Cáncer de pulmón de células
  no pequeñas; agrupa aproximadamente el 85\,\% de los casos e incluye
  adenocarcinoma, carcinoma escamoso y carcinoma de células grandes.
  \item[OMS / WHO] Organización Mundial de la Salud / World Health
  Organization. Autora del panel de referencia inmunohistoquímico
  empleado como validación externa.
  \item[PCA] Principal Component Analysis. Análisis de componentes
  principales, técnica de reducción de dimensionalidad usada en la
  visualización exploratoria por cohorte.
  \item[PD-1 / PD-L1] Programmed Death Protein 1 y su ligando. Diana
  terapéutica de la inmunoterapia con inhibidores de \emph{immune
  checkpoints}, con patrones de respuesta que dependen del subtipo
  histológico NSCLC.
  \item[PCR] Polymerase Chain Reaction. Técnica de amplificación de
  ácidos nucleicos; la variante \emph{quantitative RT-PCR} es una de
  las tecnologías candidatas para medir el panel mínimo en clínica.
  \item[RF] Random Forest. Clasificador ensamble de árboles de decisión;
  uno de los tres modelos evaluados en el trabajo.
  \item[RNA-Seq] Sequencing of RNA. Cuantificación del transcriptoma
  mediante secuenciación masiva; tecnología alternativa a los
  \emph{microarrays} y objetivo natural de la validación prospectiva.
  \item[ROC] Receiver Operating Characteristic. Curva que representa
  sensibilidad frente a $1{-}$especificidad al variar el umbral de
  decisión de un clasificador binario.
  \item[SCLC] Small Cell Lung Cancer. Cáncer de pulmón microcítico;
  representa aproximadamente el 15\,\% de los casos y queda fuera del
  ámbito del clasificador de subtipo del framework.
  \item[SQC] Squamous Cell Carcinoma. Carcinoma de células escamosas o
  carcinoma epidermoide, uno de los dos subtipos histológicos
  contemplados en la tarea de subtipo.
  \item[SVM] Support Vector Machine. Máquina de vectores soporte;
  clasificador supervisado incluido en la comparación de modelos.
  \item[TCGA] The Cancer Genome Atlas. Consorcio internacional de
  genómica del cáncer cuyos datos (LUAD y LUSC en pulmón) son
  candidatos naturales para la validación prospectiva RNA-Seq.
  \item[TFM] Trabajo Fin de Máster.
  \item[TTF-1] Thyroid Transcription Factor 1, codificado por el gen
  \gene{NKX2-1}. Uno de los marcadores IHC canónicos de
  adenocarcinoma pulmonar.
  \item[UAX] Universidad Alfonso X el Sabio, institución académica en
  la que se defiende esta memoria.
  \item[VEGF] Vascular Endothelial Growth Factor. Diana terapéutica del
  bevacizumab, cuyo uso está contraindicado en carcinoma escamoso.
\end{description}

\section*{Símbolos matemáticos}
\addcontentsline{toc}{section}{Símbolos matemáticos}

\begin{description}
  \item[$\bar{x}$] Media aritmética de una variable.
  \item[$C$] Hiperparámetro inverso de regularización en regresión
  logística. Valores menores implican mayor penalización.
  \item[$d$] $d$ de Cohen, tamaño de efecto estandarizado.
  \item[$k$] Número de folds LODO en la validación externa. En este
  trabajo $k=8$ para tumor \emph{vs.} sano y $k=3$ para subtipo.
  \item[$n$] Tamaño muestral (número de muestras en una cohorte, cohorte
  test o grupo).
  \item[$p$] $p$-valor asociado a un contraste de hipótesis. Se emplea
  ajustado por FDR (\emph{adjusted $p$-value}) en el análisis
  diferencial.
  \item[$p$ (transcritos)] Dimensión del espacio de features (número de
  genes evaluados, $\sim 22\,880$ en el análisis principal).
  \item[$\rho$] Coeficiente de correlación de Spearman. En este trabajo
  cuantifica la relación entre el score del clasificador y el contenido
  residual de pulmón sano dentro de cada tumor.
  \item[$s_p$] Desviación típica combinada (\emph{pooled}) usada en el
  cálculo del $d$ de Cohen.
  \item[$\tau$] Umbral mínimo de tamaño de efecto para que un gen entre
  en la firma consenso ($\tau = 0{,}5$ en este trabajo).
\end{description}

\section*{Etiquetas de columnas de los CSV}
\addcontentsline{toc}{section}{Etiquetas de columnas de los CSV}

\begin{description}
  \item[\texttt{Cohorte\_Test}] Identificador GEO Series
  (\texttt{GSE\ldots}) de la cohorte usada como test en un fold LODO.
  \item[\texttt{Evaluable\_Como\_Test}] Booleano que indica si una
  cohorte cumple los criterios de inclusión para actuar como test en
  LODO.
  \item[\texttt{Baseline\_Mayoritaria}] Precisión (\emph{accuracy}) del
  clasificador ingenuo que siempre predice la clase mayoritaria del
  entrenamiento.
  \item[\texttt{Ganancia\_vs\_Baseline}] Diferencia de balanced accuracy
  entre el modelo y el baseline mayoritario, en la misma cohorte test.
  \item[\texttt{En\_Panel\_Minimo}] Booleano que indica si el gen forma
  parte del panel mínimo de 20 genes.
  \item[\texttt{Marcador\_IHC\_Clinica}] Booleano que indica si el gen
  coincide con uno de los marcadores del panel IHC clínico de la OMS.
\end{description}
